% Section 20.4: Geodesics | Part: 4 | Depends on: sec:fields, sec:tangent | Status: done | Original lines: 3021-3048
%!TEX root = ../../main.tex
\subsection{Geodesics}
\begin{formal}[{geodesic}]
A curve $\gamma(t)$ is a \emph{geodesic} of $\nabla$ if $\nabla_{\dot\gamma}\dot\gamma=0$ for all $t$ (it parallel-transports its own tangent vector: ``auto-parallelism''). In coordinates:
\begin{equation}\label{eq:geodesic}
\ddot x^k+\Gamma^k{}_{ij}(x(t))\,\dot x^i\dot x^j=0 .
\end{equation}
This is a second-order ODE: given a point and an initial velocity there is a unique maximal geodesic. If one reparametrizes, $\gamma(t)=\rho(s(t))$, then $\gamma$ is a geodesic as a set iff $\rho$ satisfies $\frac{\dd^2y^k}{\dd s^2}+\Gamma^k{}_{ij}\frac{\dd y^i}{\dd s}\frac{\dd y^j}{\dd s}=f(s)\frac{\dd y^k}{\dd s}$ with $f=-s''/(s')^2$; geodesics with $f=0$ are called \emph{affinely parametrized} (and \eqref{eq:geodesic} is that case).
\end{formal}
\begin{intuition}[{the straightest possible path; free fall}]
A geodesic is a curve with zero covariant acceleration: ``no one pushes me, I just follow the geometry''. In $\R^m$ with $\nabla^0$ it is $\ddot x=0$: uniform straight motion. The term $-\Gamma^k{}_{ij}\dot x^i\dot x^j$ is the acceleration that \emph{coordinates} see; in GR it is precisely what we call the ``gravitational force'' (and in polar coordinates, the ``centrifugal'' force). In GR: \textbf{freely falling particles move along timelike geodesics; light along null geodesics.}
\end{intuition}

\begin{example}[{straight lines in polar coordinates $\Rightarrow$ centrifugal force}]\label{ex:polarstraight}
With $\Gamma^r{}_{\varphi\varphi}=-r$, $\Gamma^\varphi{}_{r\varphi}=1/r$, equation \eqref{eq:geodesic} reads
\[
\ddot r-r\dot\varphi^2=0,\qquad \ddot\varphi+\frac2r\dot r\dot\varphi=0 .
\]
The second equation gives $\frac{\dd}{\dd t}(r^2\dot\varphi)=r(\,r\ddot\varphi+2\dot r\dot\varphi)=0$: $L=r^2\dot\varphi$ is conserved (angular momentum). Then $\ddot r=L^2/r^3$: a repulsive ``centrifugal'' term appears even though nothing is pushing the particle; it is just a straight line seen in a rotating-frame-like coordinate system. \emph{Check with the line $x=d$, $y=vt$:} $r^2=d^2+v^2t^2$, $\varphi=\arctan(vt/d)$, so $r^2\dot\varphi=vd$ is constant, and $\ddot r=v^2d^2/r^3=L^2/r^3$ with $L=vd$. \checkmark
\end{example}

% REORDERED (F7): the example "geodesic equations on the unit sphere" moved to Section sec:lc (ex:spheregeod).
The geodesics of the sphere (great circles) are found in Example~\ref{ex:spheregeod} of Section~\ref{sec:lc}.
